\documentclass{article}

\usepackage{booktabs}
\usepackage{tabularx}
\usepackage{hyperref}

\title{Development Plan\\\progname}

\author{\authname}

\date{}

\input{../Comments.text}
\input{../Common.text}

\begin{document}

\maketitle

\begin{table}[hp]
\caption{Revision History} \label{TblRevisionHistory}
\begin{tabularx}{\textwidth}{llX}
\toprule
\textbf{Date} & \textbf{Developer(s)} & \textbf{Change}\\
\midrule
Date1 & Name(s) & Description of changes\\
Date2 & Name(s) & Description of changes\\
... & ... & ...\\
\bottomrule
\end{tabularx}
\end{table}

\newpage{}

This document outlines the Development Plan for \progname{} (\textbf{LingoLyrics}), a multi-platform language-learning ecosystem that bridges interactive song playback with structured language instruction. This report covers project governance, workflow and repository strategies, continuous integration pipelines, schedule breakdowns, risk management for our Proof of Concept (PoC), technology stack selections, and team operational charters.

\section{Confidential Information?}

There is no confidential information associated with this project. The project is an open capstone initiative without third-party corporate non-disclosure agreements or proprietary industry data restrictions.

\section{IP to Protect}

There is no intellectual property (IP) to protect. All source code, schema specifications, and document deliverables will be managed openly under the software engineering capstone guidelines.

\section{Copyright License}

The software, documentation, and portable file format specifications will be distributed under the \textbf{MIT License}. The open license file is included in the root of the project repository.

\section{Team Meeting Plan}

Meetings will be conducted both virtually (via Discord/Microsoft Teams) and in-person during scheduled capstone tutorial slots.

\begin{itemize}
    \item \textbf{Internal Team Meetings:} Held twice weekly. One weekly meeting will focus on sprint planning and issue assignments, while the second will serve as a mid-week sync.
    \item \textbf{Supervisor/Advisor Meetings:} Held bi-weekly or as required by the supervising professor and teaching assistants.
    \item \textbf{Meeting Structure:} Each meeting will feature a designated Chair (responsible for preparing the agenda prior to the call) and a designated Notetaker (responsible for documenting action items in GitHub Discussions/Issues).
\end{itemize}

\section{Team Communication Plan}

Primary asynchronous communication and task tracking will take place on **GitHub Issues** and **GitHub Discussions**. Day-to-day team coordination, quick queries, and ad-hoc call scheduling will take place on a dedicated team Discord server.

\section{Team Member Roles}

To maintain project velocity, roles will rotate across sprint cycles:
\begin{itemize}
    \item \textbf{Meeting Chair / Lead Facilitator:} Prepares agendas and moderates meetings.
    \item \textbf{Notetaker / Documentation Lead:} Records meeting notes and manages course report integrations.
    \item \textbf{Git / CI-CD Maintainer:} Manages branch protection rules, review enforcement, and build pipelines.
    \item \textbf{Testing \& Quality Assurance Lead:} Oversees unit test coverage, cross-platform build targets, and schema validation.
\end{itemize}

\section{Workflow Plan}

\subsection{Git and Branching Strategy}
We will use Git for version control hosted on GitHub. Our repository will adhere to a modified \textbf{GitHub Flow} branching strategy to ensure code quality and continuous integration:
\begin{itemize}
    \item \textbf{\texttt{main} Branch:} Contains stable, production-ready code. Direct commits to \texttt{main} are strictly prohibited.
    \item \textbf{\texttt{develop} Branch:} Serves as the integration branch for ongoing feature development.
    \item \textbf{Feature Branches (\texttt{feature/<issue-id>-<short-description>}):} Created off \texttt{develop} for individual features, bug fixes, or documentation tasks.
    \item \textbf{Pull Requests (PRs):} Every feature must be merged via Pull Request into \texttt{develop}. Each PR requires at least one peer review approval, passing status checks from automated testing/linting pipelines, and a resolved discussion before merging.
\end{itemize}

\subsection{Issue Management}
Project tasks, feature requests, and bug reports will be tracked using GitHub Issues:
\begin{itemize}
    \item \textbf{Issue Templates:} Standardized issue templates for Bug Reports, Feature Requests, and Task Items ensure consistent metadata collection (reproduction steps, expected behavior, acceptance criteria).
    \item \textbf{Classification and Labels:} Issues will be classified using labels categorized by:
    \begin{itemize}
        \item \textbf{Type:} \texttt{bug}, \texttt{feature}, \texttt{documentation}, \texttt{enhancement}.
        \item \textbf{Component:} \texttt{lesson-player}, \texttt{songmap-editor}, \texttt{anki-integration}, \texttt{web-repository}.
        \item \textbf{Priority:} \texttt{priority: high}, \texttt{priority: medium}, \texttt{priority: low}.
    \end{itemize}
    \item \textbf{Assignees and Milestones:} Every issue will be assigned to a specific team member and linked to a project milestone corresponding to course deliverables.
\end{itemize}

\subsection{CI/CD Pipeline}
Continuous Integration and Continuous Deployment (CI/CD) will be implemented using \textbf{GitHub Actions}:
\begin{itemize}
    \item \textbf{Continuous Integration (CI):} Triggered on pull requests to \texttt{main} and \texttt{develop}. Performs automated code linting, static code analysis, unit testing, and validates JSON/binary schema integrity for the custom song-map file format.
    \item \textbf{Automated Builds:} Cross-platform desktop and mobile client builds will be automatically compiled on tag releases to verify cross-platform compatibility across Windows, macOS, Android, and iOS.
\end{itemize}

\section{Project Decomposition and Scheduling}

\subsection{GitHub Projects}
We will use a \textbf{GitHub Projects (V2)} Kanban board to manage task progress, project velocity, and team collaboration.
\begin{itemize}
    \item \textbf{Board Views:} Tasks will be organized into standard columns (\texttt{Backlog}, \texttt{Ready for Sprint}, \texttt{In Progress}, \texttt{In Review}, \texttt{Done}).
    \item \textbf{Workflow Automation:} Cards will automatically transition across columns based on GitHub triggers (e.g., opening a PR moves an issue to \texttt{In Review}; merging a PR moves linked issues to \texttt{Done}).
    \item \textbf{Project Link:} \url{https://github.com/orgs/<Your-Org>/projects/<Project-Number>}
\end{itemize}

\subsection{Project Schedule}
The project timeline is structured into two main phases aligned with the course milestones across the Fall and Winter terms:

\subsubsection{Fall Term}
\begin{itemize}
    \item \textbf{Problem Statement \& Scope Definition:} System requirements, stakeholder analysis, and custom document selection (Design Thinking Report).
    \item \textbf{Development Plan \& Proof of Concept (PoC):} Environment setup, CI/CD pipeline configuration, core tech stack selection, and preliminary audio/lyric sync prototype.
    \item \textbf{Requirements \& Architecture Design:} Formal software architecture definition, detailed specifications for the portable song-map file format, and database schema design for the map repository.
    \item \textbf{Design Thinking Report / Interview Evaluation:} Stakeholder interviews validating learner and creator workflows.
    \item \textbf{Verification \& Validation (V\&V) Plan:} Test plan strategy covering cross-platform player usability, Anki integration testing, and file schema validation.
\end{itemize}

\subsubsection{Winter Term}
\begin{itemize}
    \item \textbf{System Architecture \& Final Design Revision:} Refinement of cross-platform player, song-map editor GUI, and language-aware text segmentation modules.
    \item \textbf{Core Feature Implementation:} Completion of lyric sync player, Anki integration handoff, song-map editor, and web repository.
    \item \textbf{V\&V Implementation \& Usability Testing:} Execution of user testing with learners and language teachers; submission of Usability Report / User Manual.
    \item \textbf{Final Capstone Demonstration \& Expo:} Final system integration, user documentation, and public presentation.
\end{itemize}

\section{Proof of Concept Demonstration Plan}

The primary risks for \progname{} include:
\begin{enumerate}
    \item \textbf{Sub-second Audio and Lyric Synchronization:} Achieving real-time, low-latency playback synchronization across non-Latin scripts (e.g., Japanese, Vietnamese) on both desktop and mobile platforms.
    \item \textbf{Language-Aware Segmentation and Transliteration:} Parsing and segmenting languages that lack spaces between words (e.g., Japanese kanji/kana tokenization) to support individual word selection and dictionary lookup.
    \item \textbf{Anki Deck Export Handoff:} Successfully transmitting selected vocabulary, context lines, and media references to external Anki instances without forcing manual user re-entry.
\end{enumerate}

To prove these risks are solvable during the Proof of Concept (PoC) demonstration, the team will show:
\begin{itemize}
    \item A working prototype playing a mapped audio track while accurately highlighting corresponding lyric lines and word tokens in real time.
    \item Successful tokenization and transliteration/romanization rendering for at least one target language with non-standard text segmentation.
    \item An automated export test generating and pushing a card with context fields into a local Anki desktop/mobile instance via AnkiConnect or file payload export.
\end{itemize}

\section{Expected Technology}

\begin{itemize}
    \item \textbf{Languages:} TypeScript (Frontend / Application Logic) and Python or Go (Web Repository Backend / Map Services).
    \item \textbf{Frontend Framework:} React Native / Expo (for cross-platform iOS and Android targets) or Electron / React (for desktop targets).
    \item \textbf{Libraries \& Utilities:}
    \begin{itemize}
        \item \texttt{kuromoji.js} / \texttt{sudachi} for Japanese text segmentation and furigana/romanization generation.
        \item Audio engine libraries for accurate timestamp tracking (e.g., \texttt{Howler.js} or native audio bindings).
    \end{itemize}
    \item \textbf{Linter Tool:} \texttt{ESLint} paired with \texttt{Prettier} for consistent code formatting across all JS/TS files.
    \item \textbf{Unit Testing Framework:} \texttt{Jest} for unit and integration testing of client components and song-map schema parsers.
    \item \textbf{CI/CD Platform:} GitHub Actions executing automated tests, linter runs, and schema checks on every pull request.
\end{itemize}

\section{Use of AI and LLMs}

\begin{itemize}
    \item \textbf{Code Development:} GitHub Copilot and LLM tools will be used for boilerplate code generation, routine utility functions, and regex creation for text parsing.
    \item \textbf{Documentation:} LLMs will be used to review text clarity, assist with LaTeX formatting, and brainstorm test cases.
    \item \textbf{Verification Protocol:} All AI-generated code must pass static analysis and unit tests, and undergo compulsory human peer review prior to merging into \texttt{develop}. No AI-generated code will be committed without explicit verification by a team member.
\end{itemize}

\section{Coding Standard}

The project will strictly follow the **Airbnb JavaScript/TypeScript Style Guide**, enforced via automated pre-commit hooks (\texttt{husky}) and GitHub Actions workflow status checks.

\newpage{}

\section{Appendix --- Generative AI Use Disclosure}

\subsection{AI Tools Used}
\begin{itemize}
    \item \textbf{Tool Name:} OpenAI ChatGPT / Claude / Google Gemini.
    \item \textbf{Version:} Current web interface models.
\end{itemize}

\subsection{How and Where Used}
\begin{itemize}
    \item \textbf{Development Plan Framing:} Assisted in structuring LaTeX subsections for workflow policies, CI/CD pipeline strategies, and project schedules based on the problem statement inputs.
    \item \textbf{Formatting:} Generated LaTeX table code and list layouts.
\end{itemize}

\subsection{Verification of AI-Generated Content}
All text and code structures generated by AI tools were thoroughly reviewed and revised by team members to guarantee alignment with McMaster capstone requirements, actual repository workflows, and project goals.

\newpage{}

\section*{Appendix --- Reflection}

\input{../Reflection.text}

\begin{enumerate}
    \item \textbf{Why is it important to create a development plan prior to starting the project?}  
    Creating a development plan establishes clear boundaries, assigns responsibilities, defines risk mitigation strategies, and ensures all team members align on project scope, tools, and timelines before writing code.

    \item \textbf{In your opinion, what are the advantages and disadvantages of using CI/CD?}  
    \textit{Advantages:} Automated feedback on broken code, early catch of syntax or logic errors, guaranteed code formatting consistency, and simplified build processes across operating systems.  
    \textit{Disadvantages:} Initial setup overhead, potential maintenance cost of pipeline scripts, and minor build queue delays on PR submissions.

    \item \textbf{What disagreements did your group have in this deliverable, if any, and how did you resolve them?}  
    Disagreements regarding tech stack selection (e.g., native vs. hybrid cross-platform desktop frameworks) were resolved by reviewing the target platform constraints specified in our goals and reaching a consensus during team meetings.
\end{enumerate}

\newpage{}

\section*{Appendix --- Team Charter}

\subsection*{External Goals}
\begin{itemize}
    \item Complete a fully functional, usable product suitable for the Capstone EXPO.
    \item Build a portfolio-ready project demonstrating cross-platform software design, language processing, and UI synchronization.
    \item Achieve high academic performance across all deliverables.
\end{itemize}

\subsection*{Attendance}

\subsubsection*{Expectations}
All members must attend scheduled internal and supervisor meetings on time. If a member anticipates being late or absent, they must notify the team at least 12 hours in advance via Discord.

\subsubsection*{Acceptable Excuse}
Medical emergencies, scheduled university exams, or documented personal events are acceptable reasons for missing meetings.

\subsubsection*{In Case of Emergency}
In emergencies, team members must inform the group as soon as reasonably possible and delegate pending tasks to another team member to avoid blocking progress.

\subsection*{Accountability and Teamwork}

\subsubsection*{Quality}
All submitted work (code and documentation) must adhere to project linting rules, pass continuous integration checks, and undergo peer review before final inclusion.

\subsubsection*{Attitude}
Team members will maintain a respectful, constructive, and open environment. Disagreements will be handled through objective technical discussions rather than personal debate.

\subsubsection*{Stay on Track}
Task completion will be reviewed weekly against GitHub Project milestone deadlines. If a task falls behind, the team will reallocate resources or adjust sprint scope to support the assignee.

\subsubsection*{Team Building}
The team will reserve dedicated informal social time before or after sprint planning meetings to build group cohesion.

\subsubsection*{Decision Making}
Technical decisions will be made by consensus. If a tie or deadlock occurs, a vote will be held after a 10-minute presentation of competing options.

\end{document}